% ═══════════════════════════════════════════════════════════════
% MINT LEHRERHINWEISE TEMPLATE v1.0
% Für Physik, Chemie, Informatik
% ═══════════════════════════════════════════════════════════════
% ENTHÄLT:
% - Kurzübersicht (Thema, Lernziele, Material)
% - Stundenverlauf (5-Minuten-Raster)
% - Differenzierung (★/★★/★★★ oder G/E-Kurs)
% - Häufige Fehler + Reaktionen
% - Material-Checkliste
% - Lösungen (Kurzform)
% - Tafelanschrieb
% - Ausblick Folgestunden
% - Reflexionsnotizen
% ═══════════════════════════════════════════════════════════════

\documentclass[11pt, a4paper]{article}

% ═══════════════════════════════════════════════════════════════
% SW-MODUS TOGGLE
% ═══════════════════════════════════════════════════════════════
\newif\ifswmodus
\swmodusfalse    % Standard: Farbe

\usepackage[utf8]{inputenc}
\usepackage[T1]{fontenc}
\usepackage[ngerman]{babel}

\usepackage{helvet}
\renewcommand{\familydefault}{\sfdefault}

\usepackage[margin=2cm, top=2cm, bottom=2cm]{geometry}
\usepackage{amsmath, amssymb}
\usepackage{siunitx}
\sisetup{locale = DE, per-mode = symbol}

\usepackage{enumitem}
\usepackage{tabularx}
\usepackage{booktabs}
\usepackage{array}
\usepackage{longtable}

\usepackage{xcolor}
\usepackage{tcolorbox}
\tcbuselibrary{skins, breakable}

\usepackage{fancyhdr}
\usepackage{lastpage}

\setlist{nosep, leftmargin=1.5em}

% ═══════════════════════════════════════════════════════════════
% FARBEN
% ═══════════════════════════════════════════════════════════════

% Fachfarben
\definecolor{physik-color}{HTML}{2c5aa0}
\definecolor{chemie-color}{HTML}{2a7a4b}
\definecolor{informatik-color}{HTML}{b35c00}

% Strukturfarben
\definecolor{hellgrau-color}{HTML}{F0F0F0}
\definecolor{grau-color}{HTML}{666666}
\definecolor{success-color}{HTML}{2a7a4b}
\definecolor{warning-color}{HTML}{b35c00}
\definecolor{error-color}{HTML}{c62828}

% Differenzierungsfarben
\definecolor{basis-color}{HTML}{2a7a4b}
\definecolor{standard-color}{HTML}{2c5aa0}
\definecolor{erweiterung-color}{HTML}{7b1fa2}

% SW-Farben
\definecolor{sw-dark}{HTML}{000000}
\definecolor{sw-gray}{HTML}{666666}
\definecolor{sw-white}{HTML}{FFFFFF}

% ═══════════════════════════════════════════════════════════════
% VARIABLEN (ANPASSEN!)
% ═══════════════════════════════════════════════════════════════

\newcommand{\fachid}{physik}                % physik | chemie | informatik
\newcommand{\fach}{Physik}
\newcommand{\thema}{Thema}
\newcommand{\klasse}{10}
\newcommand{\stunde}{01}
\newcommand{\stundenanzahl}{4}              % Gesamtzahl Stunden in Reihe
\newcommand{\reihe}{Reihenthema}
\newcommand{\dauer}{45 Min.}
\newcommand{\lerngruppengroesse}{24}
\newcommand{\datum}{\today}

% Differenzierungsmodus: sternchen | kurse | keine
\newcommand{\diffmodus}{sternchen}

% ═══════════════════════════════════════════════════════════════
% BEDINGTE FARBZUWEISUNG
% ═══════════════════════════════════════════════════════════════

\ifswmodus
    \colorlet{fachfarbe}{sw-dark}
    \colorlet{hellgrau}{sw-white}
    \colorlet{grau}{sw-gray}
    \colorlet{success}{sw-dark}
    \colorlet{warning}{sw-dark}
    \colorlet{error}{sw-dark}
    \colorlet{basis}{sw-dark}
    \colorlet{standard}{sw-dark}
    \colorlet{erweiterung}{sw-dark}
\else
    \colorlet{fachfarbe}{\fachid-color}
    \colorlet{hellgrau}{hellgrau-color}
    \colorlet{grau}{grau-color}
    \colorlet{success}{success-color}
    \colorlet{warning}{warning-color}
    \colorlet{error}{error-color}
    \colorlet{basis}{basis-color}
    \colorlet{standard}{standard-color}
    \colorlet{erweiterung}{erweiterung-color}
\fi

% ═══════════════════════════════════════════════════════════════
% KOPF-/FUSSZEILE
% ═══════════════════════════════════════════════════════════════

\pagestyle{fancy}
\fancyhf{}
\renewcommand{\headrulewidth}{0.4pt}
\lhead{\small\textcolor{fachfarbe}{\textbf{\fach{} Kl. \klasse}} -- \thema{} -- \textbf{Lehrerhinweise}}
\rhead{\small LH-\stunde}
\cfoot{\small\thepage/\pageref{LastPage}}

% ═══════════════════════════════════════════════════════════════
% BEFEHLE
% ═══════════════════════════════════════════════════════════════

% Differenzierungsmarker
\newcommand{\B}{\textcolor{basis}{\textbf{[B]}}}
\newcommand{\Std}{\textcolor{standard}{\textbf{[S]}}}
\newcommand{\E}{\textcolor{erweiterung}{\textbf{[E]}}}

% Sternchen-Marker (für Kl. 6-9)
\newcommand{\sternEins}{\textcolor{basis}{\textbf{(*)}}}
\newcommand{\sternZwei}{\textcolor{standard}{\textbf{(**)}}}
\newcommand{\sternDrei}{\textcolor{erweiterung}{\textbf{(***)}}}

% Kurs-Marker (für Kl. 9 mit Kursen)
\newcommand{\Gkurs}{\textcolor{basis}{\textbf{[G]}}}
\newcommand{\Ekurs}{\textcolor{erweiterung}{\textbf{[E]}}}

% Lösung (rot)
\newcommand{\lsg}[1]{\textcolor{error}{\textbf{#1}}}

% ═══════════════════════════════════════════════════════════════
% BOXEN
% ═══════════════════════════════════════════════════════════════

\newtcolorbox{infobox}[1][]{
    colback=hellgrau,
    colframe=fachfarbe,
    boxrule=1pt,
    arc=0pt,
    left=8pt, right=8pt, top=6pt, bottom=6pt,
    fonttitle=\bfseries,
    title=#1,
    breakable
}

\newtcolorbox{warnbox}[1][]{
    colback=white,
    colframe=warning,
    boxrule=1pt,
    arc=0pt,
    left=8pt, right=8pt, top=6pt, bottom=6pt,
    fonttitle=\bfseries,
    title=#1,
    breakable
}

\newtcolorbox{tippbox}{
    colback=white,
    colframe=success,
    boxrule=1pt,
    arc=0pt,
    left=8pt, right=8pt, top=6pt, bottom=6pt,
    fonttitle=\bfseries,
    title=Tipp,
    breakable
}

% ═══════════════════════════════════════════════════════════════
% DOKUMENT
% ═══════════════════════════════════════════════════════════════

\begin{document}

% ═══════════════════════════════════════════════════════════════
% TITEL
% ═══════════════════════════════════════════════════════════════

\begin{center}
{\Large\textbf{\textcolor{fachfarbe}{Lehrerhinweise: \thema}}}\\[0.3em]
{\normalsize \fach{} Klasse \klasse{} -- Stunde \stunde{} von \stundenanzahl{} -- Reihe ``\reihe''}
\end{center}

\vspace{1em}

% ═══════════════════════════════════════════════════════════════
% 1. KURZÜBERSICHT
% ═══════════════════════════════════════════════════════════════

\section*{1. Kurzübersicht}

\begin{infobox}
\begin{tabular}{ll}
\textbf{Thema:} & \thema{} \\
\textbf{Klasse:} & \klasse{} (Profil: heterogen) \\
\textbf{Zeit:} & \dauer \\
\textbf{Lernziele:} & [FZ1, FZ2, FZ3 aus LE] \\
\textbf{Methoden:} & [Methoden aus LE] \\
\textbf{Material:} & AB-\stunde, LP-\stunde, ggf. SIM-\stunde \\
\end{tabular}
\end{infobox}

% ═══════════════════════════════════════════════════════════════
% 2. STUNDENVERLAUF
% ═══════════════════════════════════════════════════════════════

\section*{2. Stundenverlauf (5-Minuten-Raster)}

\begin{longtable}{|c|p{2.5cm}|p{7cm}|p{3cm}|}
\hline
\textbf{Min.} & \textbf{Phase} & \textbf{Inhalt} & \textbf{Material} \\
\hline
\endhead

0--5 & Einstieg & [Beschreibung] & [Material] \\
\hline

5--15 & Input/Erarbeitung & [Beschreibung] & LP-\stunde \\
\hline

15--25 & Übung I & [Beschreibung] & AB-\stunde \\
\hline

25--35 & Übung II & [Beschreibung] & AB-\stunde \\
\hline

35--40 & Sicherung & [Beschreibung] & Tafel \\
\hline

40--45 & Abschluss & Zusammenfassung, HA & Heft \\
\hline
\end{longtable}

% ═══════════════════════════════════════════════════════════════
% 3. DIFFERENZIERUNG
% ═══════════════════════════════════════════════════════════════

\section*{3. Differenzierung}

\begin{infobox}[Legende]
\textbf{Sternchen (Kl. 6-9):} \sternEins{} Basis (BR) | \sternZwei{} Standard (MR) | \sternDrei{} Erweitert (GY)

\textbf{Kurse (Kl. 9):} \Gkurs{} G-Kurs (MR-Niveau) | \Ekurs{} E-Kurs (GY-Niveau)

\textbf{WICHTIG:} Diese Marker sind NUR für die Lehrkraft. Auf Schülermaterial (Kl. 10+) erscheinen KEINE Niveaumarkierungen!
\end{infobox}

\vspace{1em}

\begin{tabular}{|l|p{4cm}|p{4cm}|p{4cm}|}
\hline
& \sternEins{} \textbf{Basis} & \sternZwei{} \textbf{Standard} & \sternDrei{} \textbf{Erweitert} \\
\hline
\textbf{Ziel} & Qualitatives Verständnis & Formel anwenden & Formel umstellen, Transfer \\
\hline
\textbf{Aufgaben} & 1--2 (8 BE) & 1--4 (20 BE) & Alle (25 BE) \\
\hline
\textbf{Hilfen} & Formelkarte, Beispiel & Hinweise bei Umstellung & Nur Impulse \\
\hline
\textbf{LP-Niveau} & ★ (Pflicht) & ★★ (Pflicht) & ★★★ (Pflicht) \\
\hline
\end{tabular}

\vspace{1em}

\begin{tippbox}
Bei heterogenen Gruppen: Leistungsstarke SuS als ``Experten'' einsetzen, die anderen erklären können.
\end{tippbox}

% ═══════════════════════════════════════════════════════════════
% 4. HÄUFIGE SCHÜLERFEHLER
% ═══════════════════════════════════════════════════════════════

\section*{4. Häufige Schülerfehler und Reaktionen}

\begin{warnbox}[Typische Fehler]
\begin{tabular}{|p{4.5cm}|p{8cm}|}
\hline
\textbf{Fehler} & \textbf{Reaktion / Hilfe} \\
\hline
\textbf{[Fehler 1]} & ``[Reaktion/Frage zur Selbstkorrektur]'' \\
\hline
\textbf{[Fehler 2]} & ``[Reaktion/Frage zur Selbstkorrektur]'' \\
\hline
\textbf{Einheiten vergessen} & ``Jede physikalische Größe braucht eine Einheit!'' \\
\hline
\textbf{Formel falsch umgestellt} & Dreieck-Methode zeigen; ``Was steht oben allein?'' \\
\hline
\end{tabular}
\end{warnbox}

% ═══════════════════════════════════════════════════════════════
% 5. MATERIAL-CHECKLISTE
% ═══════════════════════════════════════════════════════════════

\section*{5. Material-Checkliste}

\begin{itemize}
    \item[$\Box$] AB-\stunde.pdf kopiert (\lerngruppengroesse{} Kopien)
    \item[$\Box$] ML-\stunde.pdf (1× für Lehrkraft)
    \item[$\Box$] LP-\stunde.html online/verfügbar
    \item[$\Box$] ggf. SIM-\stunde.html getestet
    \item[$\Box$] ggf. LS-\stunde.pdf (alternative zu LP)
    \item[$\Box$] Formelkarte für Niveau (★)
    \item[$\Box$] Tafelanschrieb vorbereitet
    \item[$\Box$] Experimentiermaterial (falls nötig)
\end{itemize}

% ═══════════════════════════════════════════════════════════════
% 6. LÖSUNGEN (KURZFORM)
% ═══════════════════════════════════════════════════════════════

\section*{6. Lösungen AB-\stunde{} (Kurzform)}

\textit{Vollständige Lösungen: siehe ML-\stunde.pdf}

\vspace{0.5em}

\begin{tabular}{|l|l|}
\hline
\textbf{Aufgabe} & \textbf{Lösung} \\
\hline
Merksatz & [Lücken ausgefüllt] \\
\hline
1 & [Kurzantwort] \\
\hline
2 & [Kurzantwort] \\
\hline
3 & [Ergebnis mit Einheit] \\
\hline
4 & [Kurzantwort] \\
\hline
5 & [Kurzantwort] \\
\hline
\end{tabular}

% ═══════════════════════════════════════════════════════════════
% 7. TAFELANSCHRIEB
% ═══════════════════════════════════════════════════════════════

\section*{7. Tafelanschrieb (Vorschlag)}

\begin{infobox}[Tafel links]
[Skizze/Schaltbild/Diagramm]
\end{infobox}

\begin{infobox}[Tafel Mitte]
[Formel + Erklärung]
\end{infobox}

\begin{infobox}[Tafel rechts]
[Merksatz / Zusammenfassung]
\end{infobox}

% ═══════════════════════════════════════════════════════════════
% 8. AUSBLICK FOLGESTUNDEN
% ═══════════════════════════════════════════════════════════════

\section*{8. Ausblick: Folgestunden}

\begin{tabular}{|c|l|l|}
\hline
\textbf{Std.} & \textbf{Thema} & \textbf{Schwerpunkt} \\
\hline
\stunde & \thema{} (diese Stunde) & [Schwerpunkt] \\
\hline
[\stunde+1] & [Folgestunde] & [Schwerpunkt] \\
\hline
[\stunde+2] & [Folgestunde] & [Schwerpunkt] \\
\hline
\end{tabular}

% ═══════════════════════════════════════════════════════════════
% 9. REFLEXIONSNOTIZEN
% ═══════════════════════════════════════════════════════════════

\section*{9. Reflexionsnotizen (nach Durchführung)}

\begin{tabular}{|p{7cm}|p{7cm}|}
\hline
\textbf{Was hat gut funktioniert?} & \textbf{Was muss angepasst werden?} \\
\hline
& \\[3cm]
\hline
\end{tabular}

\vspace{1em}

\textbf{Zeiteinhaltung:} $\Box$ okay \quad $\Box$ zu knapp \quad $\Box$ Zeit übrig

\textbf{Differenzierung:} $\Box$ okay \quad $\Box$ anpassen (wie: \rule{5cm}{0.4pt})

\textbf{LP/SIM:} $\Box$ funktioniert \quad $\Box$ technische Probleme \quad $\Box$ nicht genutzt

\end{document}
